\begin{textsample}{POS Dim 3 – go – Score 52.00 – go/go\_inf\_20.txt}  \label{ex:f3_pos_042}
3.4. \textbf{Escuta}, Fala, Pensamento e Imaginação Toda \textbf{criança}, desde \textbf{bebê}, tem o direito de conhecer, experienciar, compreender e se apropriar da sua língua materna, tanto no que se refere à linguagem oral, falada ou por sinais ( Língua Brasileira de Sinais - Libras ), quanto no que diz respeito à linguagem escrita, convencional ou tátil ( Braile ), ou ainda, por outros sistemas simbólicos que têm como suporte a tecnologia assistiva.
A língua como instrumento de comunicação se caracteriza como atividade humana carregada de sentidos e significados compartilhados por determinado grupo social, que possibilita a constituição do sujeito e da sua consciência, bem como o desenvolvimento do pensamento e da imaginação.
A língua é um sistema simbólico que tem como função mediar a relação dos sujeitos com o mundo e com eles mesmos.
Ela se apresenta por meio de signos, de regras, de convenções e de relações que constituem as funções psíquicas humanas, como a \textbf{percepção}, a memória, a atenção, a ação voluntária, a comunicação, o conhecimento, o raciocínio, a reflexão, a autorregulação do comportamento, a consciência e a afetividade.
Portanto, sua apropriação é um processo complexo que necessita de mediação.
A compreensão e o uso da língua pela \textbf{criança} em diferentes situações comunicativas, entendidas como espaços de negociação de significados que levem à compreensão do que se diz, possibilita que ela desenvolva a sua capacidade de simbolizar, diversificando as suas maneiras de entender e interpretar a realidade.
A BNCC ( BRASIL, 2017 ) apresenta este campo de experiências abordando conceitos centrais que permitem apreender o processo de apropriação da língua materna pela \textbf{criança}, bem como suas formas de expressar leituras, compreensões, \textbf{percepções} e sentimentos sobre o mundo, sobre os outros e sobre si mesma, conforme segue : É inerente a esse campo de experiências a discussão sobre um dos grandes desafios da Educação \textbf{Infantil}, refletir e entender qual é o papel desta etapa da Educação Básica no processo de compreensão, uso e apropriação da língua materna pela \textbf{criança}.
Dessa forma, este campo de experiência tem por objetivo abordar a articulação entre a \textbf{escuta}, a fala[19 ] e o pensamento, o desenvolvimento da imaginação, a relação da língua materna com as culturas orais e escritas e a importância da literatura no processo de constituição do sujeito. \textbf{Escuta}, Fala e Pensamento Aprender a \textbf{escutar}, a organizar o pensamento e a materializá -lo por meio da fala, \textbf{emitindo} ideias, sentimentos, \textbf{desejos}, necessidades, posicionamentos e argumentos, diante de ações, de fatos, de acontecimentos e de conhecimentos é uma atividade essencialmente humana e um dos principais aprendizados da \textbf{criança}. \textbf{Escutar}, falar e pensar têm origens diferentes no processo de aprendizagens e desenvolvimento da \textbf{criança}, mas um influencia no desenvolvimento do outro e num determinado momento passam a ser coincidentes. \textbf{Escuta} é o acolhimento e a compreensão do que o outro tem a dizer por meio de suas diferentes expressões – gestual, sonora, oral, escrita e gráfica.
Exige sensibilidade, disponibilidade e capacidade de conexão ao outro, por meio da empatia, do \textbf{desejo} e do interesse em interpretar o que está sendo dito.
Compreendida num sentido amplo, envolve o contexto em que a situação comunicativa acontece, a entonação de voz, os olhares, a postura corporal, as \textbf{vestimentas}, o local onde ela ocorre. \textbf{Escutar} pressupõe, em diferentes situações comunicativas presentes no cotidiano, desenvolver as capacidades de prestar atenção, de antecipar sentidos, de fazer conclusões, de saber respeitar a opinião e a vez do outro, de articular o pensamento e de se expressar para manter o diálogo.
A fala se constitui na transposição de \textbf{sons}, \textbf{gestos} e silêncios, em palavras, em discurso humano ( \textbf{OLIVEIRA}, 2018 ), quer seja na forma oral ou em sinais.
Aprender a usar o vocabulário de uma língua é um processo complexo, porque os \textbf{gestos}, os balbucios devem ser transformados em signos linguísticos, acordados coletivamente, que, organizados em diferentes gêneros textuais, permitem a comunicação entre os sujeitos.
De acordo com Goulart e Mata ( 2016 ), a \textbf{criança} quando chega ao mundo, antes mesmo de nascer, já é referida pelos familiares e pelas pessoas que convivem com a mãe quando \textbf{conversam} com ela ainda na barriga.
Ou seja, desde a concepção, a \textbf{criança} já faz parte de um grupo social, que possui determinados \textbf{hábitos}, costumes e uma língua.
Assim, é nas interações cotidianas com essa língua, que se materializa por meio da palavra, compreendida como meio de comunicação e como um conceito e uma generalização ( VIGOTSKI, 2000 ), que a \textbf{criança} aprende a falar e que o mundo é apresentado, nomeado e significado para e por ela.
Segundo o autor citado, a palavra constitui a consciência e é o elo entre a fala e o pensamento.
O pensamento constitui -se numa ação do intelecto, que possibilita representar mentalmente conhecimentos, fatos, ideias, opiniões, sentimentos, \textbf{desejos}, necessidades, interesses e emoções.
Sua materialização ocorre por meio da palavra, de forma que pensamento e palavra possuem uma relação indissociável e interdependente, mesmo que, em sua origem, os seus processos não sejam coincidentes.
Essa inter-relação entre pensamento e palavra se dá por meio dos significados, que são conceitos estabelecidos socialmente e que permitem a comunicação e o diálogo entre as pessoas.
No entanto, os significados modificam -se ao longo do processo de aprendizagens e desenvolvimento de cada sujeito.
Para uma \textbf{criança} de dois anos o significado de uma palavra é bem diferente do que quando maior ou na fase \textbf{adulta}.
Pode -se citar, por exemplo, a palavra moradia que significa local onde as pessoas habitam para se proteger do sol, do frio ou da chuva.
Para uma \textbf{criança} \textbf{pequena} suas referências serão as casas das pessoas com quem convive, conforme cresce esse conceito é modificado, à medida que estuda sobre os diferentes tipos de moradia da sua época e de outras, conhecendo e visualizando variadas imagens.
Isso ocorre porque os significados são construídos a partir das experiências vivenciadas pelos sujeitos na interação com o outro e com o que significam as palavras em situações concretas.
Como a \textbf{criança} aprende a \textbf{escutar}, a falar e a pensar?
De que forma esses processos são desenvolvidos?
A princípio, a \textbf{criança} se manifesta por meio de choros, balbucios, risadas, gritinhos e expressões faciais.
Essas são suas formas de interagir com o mundo, na busca de apreendê -lo.
Dessa forma, a \textbf{criança} percebe o mundo por meio do olhar, dos movimentos corporais, do toque, do cheiro, da experimentação de diferentes sabores, \textbf{texturas}, \textbf{percepção} dos \textbf{sons} e dos silêncios, do diálogo e da interação com o outro.
Mas, só perceber o mundo não é suficiente, é preciso que alguém o interprete para que ela possa significá -lo.
Transformar o sentir, o perceber, o explorar em linguagem, portanto em significados e sentidos, é o que torna o mundo \textbf{compreensível} para a \textbf{criança}.
Como a \textbf{criança} decifra e interpreta o mundo?
Nas relações sociais.
Desde que nasce ela se empenha em compreender os signos utilizados pelos humanos ao seu redor e procura apropriar -se deles para se fazer entendida pelos outros.
Ao chorar, a mãe, ou o outro \textbf{adulto} que \textbf{cuida} da \textbf{criança}, interpreta esse choro como desconforto dizendo “ Está com fome?
Quer papá, querido? ” ou “ Quer colinho?
Quer sair desse berço?
Venha com a mamãe! ”.
Ao realizar essa ação, transforma em palavras esse incômodo da \textbf{criança} e dá sentido às suas expressões, por meio desse \textbf{gesto} de empatia.
Ou seja, essas expressões ganham intencionalidade a partir das interpretações e reações dos outros perante elas, que, além de traduzir em palavras as sensações do \textbf{bebê}, pode ainda oferecer seu \textbf{colo}, seu carinho, seu olhar, sua raiva e irritação.
Com o tempo e as constantes significações que o \textbf{adulto} vai dando aos \textbf{gestos}, expressões faciais, choros dos \textbf{bebês} no \textbf{dia} a \textbf{dia}, eles vão aprendendo a interpretar suas próprias sensações corporais elaborando representações sobre as situações vivenciadas e desenvolvendo a capacidade de pensar e de se expressar.
Assim, aprender a acalmar, a chorar mais alto, a esperar, a sentir -se querido ou rejeitado, a balbuciar, a produzir as primeiras palavras com significados difusos relativos aos contextos situacionais - qué, dá, bô -, a combinar palavras iniciando o processo de gramaticalização, flexão e sintaxe - qué aua, dá neném, fazi cocô - e, gradativamente, descobrir que tudo tem um nome e que há uma certa estabilidade nos significados das palavras é muito complexo para a \textbf{criança} e, a isso, se denomina o desenvolvimento da fala.
Vygotsky ( 2000 ), ao abordar o desenvolvimento do pensamento e da fala na \textbf{criança}, esclarece que a \textbf{criança} no início do processo de verbalização tem a necessidade de falar enquanto age.
Isto é, enquanto a \textbf{criança} realiza algo, ela diz para si mesma o que está fazendo.
Vygotsky denomina essa fala como organizadora do pensamento e da ação da \textbf{criança}.
À medida que planeja mentalmente suas ações, a \textbf{criança} elabora o seu discurso interior, que se configura como o início do pensamento de forma mais elaborada.
Começa, então, a verbalizar para o outro o que faz.
O autor citado denomina esse processo de fala comunicativa.
Ao distinguir a fala para si da fala para o outro, a \textbf{criança} materializa o pensamento na palavra e, ao mesmo tempo, pode modificar o seu pensamento ao verbalizá -lo.
Assim, apropriar -se da linguagem oral é mais do que apreender um idioma, é apropriar -se de formas de expressão e comunicação, formas estas que constituem o pensamento e a imaginação.
A \textbf{criança} amplia seus conhecimentos, desenvolve novas estruturas linguísticas e aprende a construir narrativas a partir da qualidade, da quantidade e da variedade de situações comunicativas com as quais ela convive.
Nesse processo, passa das nomeações para as descrições de fatos e de situações até \textbf{conseguir} construir narrativas mais elaboradas com uma sequência lógica e autoral.
A \textbf{criança} surda se apropria do significado de expressões, objetos, situações, fatos, acontecimentos, por meio da palavra proferida pelos \textbf{gestos}.
É importante ressaltar que, uma vez detectada a surdez, a \textbf{criança} e a sua família têm o direito de se apropriarem da Libras.
Quanto mais cedo essa língua for inserida no processo de aprendizagens e desenvolvimento da \textbf{criança}, mais possibilidades de significação do mundo ela terá.
A ampliação do vocabulário e o domínio progressivo do uso das palavras em situações comunicativas, utilizando estruturas linguísticas cada vez mais elaboradas, se dará com a promoção de interações com as diferentes linguagens.
Isto é, a \textbf{criança} produz e se apropria da língua quando chamada a falar, a dizer ou contradizer, a escutar/ouvir, a descrever, a contar/narrar, a opinar, a concordar ou discordar, a argumentar, a propor, a definir, a compreender e participar efetivamente do mundo.
Na relação com todas as \textbf{crianças}, sejam surdas, cegas, com paralisia cerebral, transtorno ou síndrome que interfere na produção e apropriação da língua materna, o desafio é construir formas de compreender suas emoções, suas formas de significar e interpretar o mundo e de se fazer compreendido por elas, a fim de garantir o desenvolvimento da linguagem, do pensamento e da imaginação.
Imaginação O que é imaginação?
Qual sua importância para o processo de aprendizagens e desenvolvimento?
Diferente do que se costuma acreditar, a imaginação não é, necessariamente, sinônimo de fantasia, no sentido de pensá -la como algo fora da realidade.
A fantasia se constitui de uma combinação de elementos retirados da realidade vivida que são transformados e reelaborados pela imaginação.
Isso significa que a imaginação é um processo de invenção, de descobertas e de criação que, \textbf{apoiado} na realidade e na produção cultural, ou seja, no que foi vivido anteriormente pelo sujeito, permite a ele fazer recombinações e (re)elaborações produzindo algo novo, quer seja um objeto externo ou uma construção da mente ou da emoção ( VIGOTSKI, 2009 ).
A imaginação é fundamental para a sobrevivência da espécie humana, porque sem ela não seria possível a construção do mundo sociocultural, uma vez que tudo que não pertence ao mundo natural consiste em produção humana desenvolvida a partir da sua atividade criadora.
Por isto, só existem casas, \textbf{carros}, aviões, submarinos, tanques de guerra, armas etc., porque alguém teve a ideia e usou a imaginação para criá -los.
Diante da sua importância, o desenvolvimento da imaginação deve ser uma ação valorizada e garantida nas \textbf{instituições} de Educação \textbf{Infantil}.
Isso se dará à medida que as ações pedagógicas forem as mais variadas possíveis, envolvendo as diferentes linguagens e os conhecimentos do patrimônio da humanidade, assim como a \textbf{criança} ter liberdade para experimentar, \textbf{manusear} diferentes materiais, \textbf{texturas}, (re)organizar espaços e objetos, expressar suas ideias, sendo autorais em seus posicionamentos e em suas produções, quer sejam textos orais, desenhos, escritas etc.
Nessa perspectiva, para a \textbf{criança} imaginar e criar é necessário conhecer, significar, \textbf{imitar}, sentir, cheirar, provar, apalpar, se \textbf{emocionar}.
Desse modo, se alguém propuser a uma \textbf{criança} que \textbf{imite} um cachorro, ela terá que evocar na memória imagens, cheiros, \textbf{sons} \textbf{emitidos}, filmes assistidos e vivências que teve com esse animal, para que possa ter uma imagem mental de um cachorro.
Logo, terá que ter conhecido um antes, senão pessoalmente, por meio de representações imagéticas.
Quanto mais experiências tiver com cachorro, mais saberá que é um animal que se apresenta de diferentes formas, pequeno-grande, calmo-agitado, manso-bravo, peludo-pouco pelo.
Assim, a \textbf{criança} terá à sua disposição diferentes elementos para realizar a imitação.
O mesmo ocorre no processo de representação gráfica do cachorro, no desenho, na \textbf{pintura}, em telas, computadores, tablets, papéis etc.
Nas \textbf{brincadeiras} de faz de conta, a imaginação e a fantasia expressam a capacidade da \textbf{criança} de \textbf{imitar} o real, mas não da forma como foi vivenciado, e sim, com vários elementos diferentes, resultantes da sua criação.
Esse processo, de imaginar e de fantasiar, auxilia a \textbf{criança} a compreender a realidade e os diferentes papéis que os sujeitos ocupam nos grupos sociais.
Assim, a \textbf{criança}, ao recordar e reviver suas experiências em diferentes contextos, cria novas combinações e possibilidades de interpretar e representar o real, de acordo com seus \textbf{desejos}, suas necessidades, seus \textbf{afetos} e suas emoções.
Em síntese, a matéria-prima para a imaginação criadora é a memória das experiências vividas proporcionada pela participação efetiva da \textbf{criança} em práticas sociais.
Língua Materna e suas Relações com as Culturas Orais e Escritas A apropriação e a compreensão da língua materna, como já mencionado, é um processo complexo e acontece por meio das interações que ocorrem nas práticas sociais, a partir da utilização da linguagem verbal, que corresponde à linguagem oral e à linguagem escrita.
Essas duas linguagens se inter-relacionam, mas não são iguais, possuem características próprias e saberes específicos que as definem.
Ambas são importantes e necessárias para a construção cultural de um povo e para a organização da vida coletiva.
Contudo, dependendo da constituição dos grupos sociais uma linguagem pode ser mais valorizada e utilizada que a outra.
Em determinadas comunidades basta a palavra proferida oralmente para ser realizada uma transação comercial, em outras, é necessário contrato com assinatura para que o negócio se efetive.
Esses diferentes usos da linguagem verbal não podem levar à consideração de que uma linguagem é superior à outra.
É preciso ressaltar esse aspecto, porque, em sociedades em que a escrita é utilizada de forma sistemática há uma tendência em não valorizar a linguagem oral.
Na organização da ação pedagógica das \textbf{instituições} de Educação \textbf{Infantil}, a linguagem verbal deve ser trabalhada com igual importância no processo de aprendizagens e desenvolvimento das \textbf{crianças}, por meio das culturas orais e das culturas escritas.
Mas, o que elas significam?
O que diferencia uma da outra?
Culturas Orais As culturas orais consistem nos saberes, fazeres e conhecimentos produzidos pelos grupos sociais e que são repassados de geração a geração por meio da linguagem oral, compondo sua memória e identidade.
Nos processos de socialização e apropriação desse repertório cultural, são utilizados causos, fábulas, parlendas, provérbios, contos, versos, parábolas, canções, mitos, rezas, representações em folguedos, encenações, como uma forma de entender a realidade, seus fenômenos e acontecimentos.
Esses gêneros textuais, geralmente, são repassados \textbf{boca} a \textbf{boca} e não há um autor em específico, porque à medida que o outro se apropria, elementos são retirados ou acrescentados de acordo com sua realidade, tendo muitas variações para um mesmo texto falado.
Assim, autoria é denominada como de “ domínio popular ”, uma vez que pertence à memória e à cultura de um grupo.
Uma das características da linguagem oral é a variedade linguística utilizada pelos sujeitos nas situações comunicativas, devido à diversidade de contextos, de intenções e de usos da língua, na maioria dos casos, a proximidade dos interlocutores pela presença física do outro, possibilitando esclarecer melhor ou \textbf{repetir} o que foi dito, perceber os \textbf{gestos} de quem está falando e a entonação da voz - \textbf{sugestão}, pedido, ordem, ameaça, conversa -, bem como estabelecer uma referência comum para o que está sendo dito.
Diferente da linguagem escrita, em que o texto tem que ser o mais claro possível para que o outro compreenda o significado, porque o leitor não terá a presença física do autor.
A linguagem oral é composta por situações comunicativas informais e também por situações públicas de comunicação oral, em que há \textbf{registros} que a precedem e o monitoramento da fala enquanto é proferida, como em palestras, conferências, mesas redondas, debates, apresentações temáticas, seminários, programas de rádio, de televisão, do YouTube.
Considerar a linguagem oral nessa totalidade significa reafirmar que, tanto em situações comunicativas formais quanto informais, há intencionalidade no dizer, não sendo, portanto, uma ação espontânea.
Na linguagem oral, o \textbf{escutar}, o falar, o \textbf{repetir}, o \textbf{imitar} e o observar são fundamentais, para que os sujeitos sejam capazes de se comunicar.
Nas \textbf{instituições} de Educação \textbf{Infantil} essas aprendizagens podem ser promovidas quando a \textbf{criança} é incentivada a : dialogar com outra ; se expressar nas \textbf{rodas} de conversas ; socializar com outras turmas o que foi aprendido ; transmitir recados ; recontar histórias ; recitar parlendas, versos, poemas, trava-línguas ; \textbf{brincar} de faz de conta, em que ela tem que se colocar diante do outro - posto de gasolina, táxi, pet shop, \textbf{feiras}, caixa de supermercado, médico - ; elaborar, com auxílio do(a) professor, as notícias faladas da turma ou da \textbf{instituição} para serem socializadas com as famílias ; entre tantas outras situações.
A \textbf{criança} que tem a possibilidade de expressar \textbf{desejos} e sentimentos, de desenvolver sua capacidade de argumentação, de dar respostas, de questionar, de avaliar fatos, situações, sem reprimendas dos \textbf{adultos}, desenvolve \textbf{autoconfiança} e maior segurança para se posicionar diante dos outros.
Culturas Escritas Cultura escrita, de acordo com Galvão ( 2016 ), se refere ao lugar simbólico e material que a escrita ocupa em determinado grupo social.
Ou seja, está relacionada à importância e ao valor atribuído pelos sujeitos à escrita, bem como à sua utilização no \textbf{dia} a \textbf{dia}.
A escrita é valorizada de forma diferente pelas pessoas, a depender dos costumes, valores, \textbf{hábitos} e organização dos grupos sociais aos quais elas pertencem, sendo necessário utilizar a terminologia no plural, “ culturas escritas ”, para contemplar essa diversidade.
Por que é importante para o professor compreender o que são as culturas escritas e no que elas influenciam no desenvolvimento da ação pedagógica com as \textbf{crianças}?
A \textbf{criança}, desde \textbf{bebê}, participa de práticas sociais, em que a leitura e a escrita se fazem presentes, porque vive em uma sociedade organizada por e a partir da linguagem escrita, denominada grafocêntrica.
No entanto, como já foi apontado, a forma como essa participação se efetiva está relacionada aos modos como sua comunidade utiliza a escrita.
O que lê e escrevem?
Quando e para que usar a leitura e a escrita?
Quais gêneros textuais são mais comuns e com quais suportes opera?
As respostas a essas questões indicam que em sociedades complexas e heterogêneas, o espaço do escrito envolve relações de poder, é múltiplo, desigual e se transforma permanentemente.
As \textbf{instituições} de Educação \textbf{Infantil} têm um papel fundamental no processo de possibilitar à \textbf{criança} o conhecimento e a participação em diferentes práticas sociais em que o escrito tem uma função real, buscando equiparar suas oportunidades, assim como entender o significado e a importância que o escrito ocupa no \textbf{dia} a \textbf{dia} dela, garantindo a compreensão das culturas escritas, para que, desde \textbf{pequena} se aproprie do como se escreve, para que e para quem se escreve.
Nessa perspectiva, qual é o papel da Educação \textbf{Infantil} no processo de compreensão e apropriação da linguagem escrita?
Nas DCNEI/2009 ( Parecer CEB n. 20 ) é afirmado que a \textbf{criança} tem interesse pela linguagem escrita, antes mesmo do trabalho formal realizado pela \textbf{instituição} de Educação \textbf{Infantil} e apresenta que > Sua apropriação pela \textbf{criança} se faz no reconhecimento, compreensão e fruição da linguagem que se usa para escrever, mediada pela \textbf{professora} e pelo professor, fazendo -se presente em atividades prazerosas de contato com diferentes gêneros escritos, como a leitura \textbf{diária} de \textbf{livros} pelo professor, a possibilidade da \textbf{criança} desde cedo \textbf{manusear} \textbf{livros} e revistas e produzir narrativas e ‘ textos ’, mesmo sem saber ler e escrever. ( p.94 ) Diante dessas indicações, apresentadas num documento de caráter normativo, o papel da Educação \textbf{Infantil} no processo de apropriação e de compreensão da linguagem escrita é de : - propiciar a participação das \textbf{crianças} em práticas sociais de leitura e de escrita de forma a favorecer a compreensão dos seus usos e funções na sociedade ; - proporcionar às \textbf{crianças}, desde \textbf{bebês}, a apreciação e a convivência \textbf{diária} e sistemática com a contação e a leitura de histórias, ampliando seu repertório literário e o conhecimento de autores e ilustradores variados ; - favorecer por meio da organização de tempos, espaços e materiais que as \textbf{crianças} explorem e \textbf{brinquem}, utilizando objetos que fazem parte das culturas escritas, como \textbf{embalagens}, cadernetas, \textbf{livros}, agendas, \textbf{lápis}, canetas, revistas, jornais, folders, caixas, bulas de remédios, cadernos ; - possibilitar às \textbf{crianças}, mediadas pelo professor, a leitura e a escrita de forma espontânea[20 ] e convencional, de gêneros textuais variados, em que elas possam compreender sua estrutura, bem como os diferentes suportes que podem ser veiculados, inclusive o digital, de acordo com a intenção comunicativa ; - aguçar o interesse e a \textbf{curiosidade} das \textbf{crianças} pela leitura e pela escrita, por meio de situações comunicativas reais, em que elas possam iniciar a compreensão das regras do Sistema de Escrita Alfabético, como por exemplo, que se escreve de cima para baixo, da esquerda para a direita, que são utilizadas letras para escrever e não outros símbolos etc.
Portanto, é necessário reafirmar que a Educação \textbf{Infantil} tem o papel de oportunizar à \textbf{criança} a compreensão dos diversos usos e funções que a escrita tem na sociedade, a \textbf{percepção} de regularidades, como as citadas e a estrutura formada por letras e sinais gráficos.
Cabe ainda ressaltar que a linguagem escrita deve ser trabalhada com a mesma importância das demais linguagens – visuais, corporal, gráfica, musical, audiovisual etc. – em contextos reais e significativos.
A tarefa do Ensino Fundamental é a alfabetização stricto sensu, sendo compreendida como processo de apropriação, compreensão e efetivo uso do complexo Sistema de Escrita Alfabético a partir de um trabalho pedagógico explícito e sistemático no entendimento do seu funcionamento, suas regularidades e irregularidades.
Diante da definição do papel da Educação \textbf{Infantil}, cabe perguntar : como as \textbf{crianças} participam de práticas de leitura e de escrita? \textbf{Crianças} \textbf{pequenas} lêem e escrevem, o que e como lêem e escrevem?
Todas as \textbf{crianças} lêem o mundo desde o ventre da mãe.
Lêem vozes, \textbf{sons}, expressões faciais, cheiros, temperatura do ambiente, humor de quem \textbf{cuida} delas, \textbf{colo}, \textbf{gestos} de carinho, afago, irritação e rispidez.
À medida que essas leituras vão sendo feitas, elas buscam interpretá -las e vão atribuindo sentidos ao que lhes acontecem.
Quando, no momento da troca, o professor se coloca à disposição da \textbf{criança}, mantém contato visual, \textbf{acolhe} os \textbf{gestos} que ela faz de estender o braço e tentar tocar o seu rosto, e esse se aproxima, confirmando o toque que ela busca e conversa com a \textbf{criança}, ele estabelece um diálogo e possibilita a descoberta da linguagem.
Todas as expressões, visuais, táteis ou sonoras, que forem ofertadas às \textbf{crianças} funcionam como ‘ \textbf{livro} ’, material a ser lido e interpretado pela \textbf{criança} e pelo \textbf{adulto}.
Assim, compreende -se leitura num sentido amplo, em que a leitura de mundo, como já afirmava Paulo Freire, precede a leitura da palavra, sendo a leitura relacionada à compreensão dos signos linguísticos, extensão e consequência, dessas diferentes leituras.
Nesse sentido, é importante ressaltar que outros gêneros textuais e suportes têm surgido, com o uso cada vez mais \textbf{frequente} das tecnologias digitais no \textbf{dia} a \textbf{dia}, exigindo novas formas de leitura.
É função das \textbf{instituições} de Educação \textbf{Infantil} e dos ( as ) professores ( as ) proporcionarem à \textbf{criança}, desde \textbf{bebê}, que participe de variadas práticas de leituras e de \textbf{registros}, em que tenha possibilidade de expressar opiniões sobre o que gostou ou não, qual \textbf{personagem} gostaria de ser, se fosse o autor qual final seria dado à história, de \textbf{brincar} com palavras por meio de rimas e aliterações, de \textbf{conversar} sobre quais estratégias utiliza para ganhar em jogos eletrônicos, entre outras.
De acordo com Vygotsky ( 2000 ), a escrita por ser uma representação simbólica não pode ser compreendida de forma dissociada de outras linguagens.
O seu processo de apropriação tem origem muito antes da \textbf{criança} \textbf{conseguir} escrever de forma convencional utilizando letras, porque a produção de textos por meio de signos inicia -se pela \textbf{criança} a partir da utilização do próprio corpo com \textbf{gestos} \textbf{indicativos} e expressões corporais que são interpretadas pelo outro e que posteriormente, são incorporados aos desenhos e às \textbf{brincadeiras} de faz de conta.
Essas diferentes representações utilizadas pela \textbf{criança} ( \textbf{gestos}, desenho, faz de conta ), somadas às vivências com \textbf{pinturas}, colagens, modelagens, utilização de letras e outros símbolos, imagens e situações vividas, observadas ou imaginadas lhes possibilitam criar textos orais que podem ser escritos por ela de forma espontânea ou ditados para um escriba transcrever.
De acordo com essa perspectiva, a \textbf{criança}, desde \textbf{bebê}, é capaz de produzir textos, de deixar \textbf{marcas}, portanto, lê e escreve de formas variadas, até o momento em que suas \textbf{marcas} se constituem em signos linguísticos acordados socialmente expressos na forma de uma escrita convencional.
Para tanto, é preciso que a \textbf{criança} perceba a importância e a necessidade de aprender a escrever, sendo fundamental que ela participe de situações comunicativas em que a linguagem escrita é utilizada, como : escrita de bilhetes para as famílias ; identificação de \textbf{brinquedos}, de objetos da sala e de pertences pessoais ; produção de cartazes para socializar com outras \textbf{crianças} o que está sendo aprendido ; lista de histórias favoritas ; produção de poemas para recitar num sarau ; produção de \textbf{livro} da turma com recontos de histórias ; entre outras.
Literatura A literatura se apresenta como possibilidade do leitor vivenciar a alteridade, deslocar -se para viver outras vidas, outros mundos imaginários novos e surpreendentes.
Possibilita, assim, um conhecimento mais profundo do humano ao aproximar o leitor dos sentimentos, das emoções, das situações, dos questionamentos, das inquietações trazidas pelo outro ( \textbf{personagem}, contexto ), enriquecendo sua própria experiência.
Pela capacidade de potencializar processos de \textbf{humanização}, a literatura será tratada nesse documento como direito inalienável do sujeito ao acesso, fruição e apropriação.
Ítalo Calvino ( 1990 ) aponta algumas características da literatura que auxilia o leitor, nesse caso a \textbf{criança}, no processo de \textbf{humanização} : - a leveza do texto literário é a possibilidade de tratar de temas complexos e difíceis como a violência, a morte, a separação, o abandono, o desprezo, de forma artística, colocando a \textbf{criança} na dimensão do sensível, permitindo que ela elabore o assunto a partir de uma outra lógica ; - a exatidão é a capacidade que o texto literário tem de organizar as palavras de forma concisa e precisa para expressar pensamentos, raciocínios, ideias, emoções ; - a rapidez está vinculada ao tempo literário, que não é cronológico, porque pode se passar muitos anos de um parágrafo para outro ou pode imputar à \textbf{criança} a espera para saber o que vai acontecer.
Apresenta, também, sucessão, ordem, simultaneidade ; - a visibilidade é a característica que oportuniza a imaginação do narrado, a criação de cenários e a corporificação de \textbf{personagens}, despertando -lhe sentimentos e sensações reais.
Também inclui a relação entre o texto e a ilustração que provoca leituras dos não ditos e novas interpretações ; - a multiplicidade revela os mistérios presentes na literatura como ser um e muitos ao mesmo tempo, ser a princesa ou o príncipe e ser a bruxa malvada.
Viver várias vidas se tornando outros, enriquecendo a própria existência ; - a consistência está relacionada à qualidade do que é oferecido ao leitor, as escolhas dos temas, dos \textbf{livros} e das formas como serão tratados.
Nesse sentido, pensar o trabalho com a literatura nas \textbf{instituições} de Educação \textbf{Infantil} implica superar uma visão instrucional e pragmática, na qual lê -se para ensinar as letras ou para a \textbf{criança} aprender a se comportar.
Compreendê -la como lugar : de relações ; de \textbf{brincadeiras} ; de produção de sentido e significados compartilhados ; de conhecimento de si e do outro ; de constituição da subjetividade e de sentimento de pertença ; de ampliação das experiências com a cultura escrita ; é a possibilidade de efetivar um trabalho diferenciado com a literatura, na perspectiva de garantir esse direito inalienável do sujeito, de se tornar mais humano por meio de escritos e imagens de outrem.
A literatura é capaz de deslocar o leitor/ouvinte quando provoca emoções, risos e choros, quando promove o diálogo entre o leitor/ouvinte e a obra por meio das ilustrações, da forma como o texto é organizado no \textbf{livro}, das expressões, metáforas, rimas e aliterações usadas pelo autor, provocando efeitos de sentidos.
Em síntese, a leitura literária não é adorno, nem entretenimento para acalmar e divertir as \textbf{crianças}, muito menos um servir de pretexto ou instrumento de ensino da língua escrita ou de outros conteúdos.
Ela possibilita o conhecimento de si e do outro, a ampliação do que já sabe do/sobre o mundo, a organização de experiências pela palavra, a imaginação, a criação, o pensamento, a emoção, a apreciação estética dos textos verbais e visuais e suas interlocuções. [ 19 ] : O uso dos termos “ fala ” ou “ palavra ” neste documento se refere à expressão do sujeito por meio da linguagem oral ou da Língua Brasileira de Sinais. [ 20 ] : A escrita espontânea se refere aos \textbf{registros} gráficos das \textbf{crianças} – rabiscos, garatujas, bolinhas, símbolo, números, letras – realizados de forma mais livre e com a intenção de escrever.
Nesse processo, elas mobilizam seus conhecimentos para compreender como se escreve, elaborando ideias, teorias a respeito do funcionamento da escrita alfabética.

% matched lemmas: acolher, adulto, afeto, apoiar, autoconfiança, bebê, boca, brincadeira, brincar, brinquedo, carro, colo, compreensível, conseguir, conversar, criança, cuidar, curiosidade, desejo, dia, diário, embalagem, emitir, emocionar, escuta, escutar, feira, frequente, gesto, humanização, hábito, imitar, indicativo, infantil, instituição, livro, lápis, manusear, marca, oliveira, pequeno, percepção, personagem, pintura, professora, registro, repetir, roda, som, sugestão, textura, vestimenta
\end{textsample}
